\documentclass[10pt]{article}

\usepackage[utf8]{inputenc}

\usepackage[T1]{fontenc}

\usepackage{amsmath}

\usepackage{amsfonts}

\usepackage{amssymb}

\usepackage[version=4]{mhchem}

\usepackage{stmaryrd}

\usepackage{graphicx}

\usepackage[export]{adjustbox}

\graphicspath{ {./images/} }



\begin{document}

то в качестве штрафюой функции можно взять



\[

\begin{gathered}

P(u)=\int_{t_{0}}^{T}\left(\sum _ { i = 1 } ^ { m } \left(\max \left\{g_{i}(x(t ; u), t) ; 0\right)^{p}+\right.\right. \\

\left.+\sum_{i=m+1}^{s}\left|g_{i}(x(t ; u), t)\right|^{p}\right) d t, p \geqslant 1,

\end{gathered}

\]



и задачу (1)-(3), (9) заменить задачей минимизации функционала $\Phi_{k}(u)=J(u)+A_{k} P(u)$ при условиях (2), (3), где $A_{k}$ - птрафной коэффициент, $\lim A_{k}=+\infty$.



Другие методы решения задачи (1) $\left.\stackrel{k}{ } \overrightarrow{3}^{\infty}\right)^{\infty},(9)$ основаны на принципе максимума Понтрягина, на динамич. программировании (см. Понтрягина принцип максияума, Динамическое программирование, Вариационное исчисление; численные методы).



Для решения задачи минимизаций квадратичного функциовала на ретениях систем линейных обыкновенных дифференциальных уравнений или линейных уравнений с частными производными, может быть применен метод моментов (см. [3], [8]). Ниже описан этот метод применительно к задаче минимизации функционала



\[

J(u)=|x(T ; u)-y|_{E^{n}}^{2},

\]



где $x=x(t ; u)$ - решение задачи $\dot{x}=A(t) x+B(t) u(t)+f(t), t_{0} \leqslant t \leqslant T, x\left(t_{0}\right)=x_{0}$,



управления $u=u(t) \in L_{2}^{r}\left[t_{0}, T\right]$ таковы, что



\[

\|u\|_{L_{2}^{r}}^{2}=\int_{t_{0}}^{T}|u(t)|_{E^{r}}^{2} d t \leqslant R^{2}

\]



здесь $A(t), B(t), f(t)$-заданные матрицы порядка $n \times n, n \times r, n \times 1$ соответственно, имеющие кусочно непрерывные әлементы на отрезке $t_{0} \leqslant t \leqslant T ; 0<R \leqslant \infty$,



\includegraphics[max width=\textwidth, center]{2023_07_09_1ad6f38ee0d6afef83f1g-1}

$=\sum_{i=1}^{m} x^{i} y^{i}$-скалярное шроизведение в $E^{m}$. Из шравнла множителей Лагранжа следует, что управлевие $u=u(t)$ является оптимальным в задаче (10)-(12) тогда и только тогда, когда существует число $\gamma \geqslant 0$ (Лагранжа множитель для ограничения (12)) такое, что



\[

\begin{gathered}

J^{\prime}(u)+\gamma u=0, \\

\psi\left(\|u\|_{L_{2}^{r}}-R\right)=0, \quad\|u\|_{L_{2}^{r}} \leqslant R ;

\end{gathered}

\]



при $R=\infty$ здесь $\gamma=0$. Из формул (5)-(8) следует, тто градиент $J^{\prime}(u)$ функционала (10) в $L_{2}^{r}\left[t_{0}, T\right]$ имеет вид $J^{\prime}(u)=B^{T}(t) \psi(t ; u), t_{0} \leqslant t \leqslant T$,



где $\psi=\psi(t ; u)$ - решение задачи



\[

\begin{gathered}

\dot{\psi}=-A^{T}(t) \psi, t_{0} \leqslant t \leqslant T, \\

\psi(T)=2(x(T ; u)-y) ;

\end{gathered}

\]



$A^{T}, B^{T}$ - матрицы, полученные транспонированием матриц $A, B$ соответственно. Условие (13) тогда примет вид



\[

B^{T}(t) \psi(t ; u)+\gamma u(t)=0, t_{0} \leqslant t \leqslant T .

\]



Условие (16) равносильно соотношениям



$\int_{t_{0}}^{T}\left\langle u(t), B^{T}(t) p_{k}(t)\right\rangle_{E^{r}} d t-\frac{1}{2}\left\langle\psi(T), p_{k}(T)_{E^{n}}=a_{k}\right.$,



\[

k=1, \ldots, n \text {, }

\]



где



\[

\begin{gathered}

a_{k}=\left\langle y, p_{k}(T)\right\rangle_{E^{n}}-\left\langle x_{0}, p_{k}\left(t_{0}\right)\right\rangle_{E^{n}}- \\

-\int_{t_{0}}^{T}\left\langle f(t), p_{k}(t)\right\rangle_{E^{r}} d t,

\end{gathered}

\]



$p_{k}(t)$ - решение системы (15) при условии $p_{k}(T)=$ $\stackrel{e_{k}}{ }=(0, \ldots, 0,1,0, \ldots, 0)$ - единичный вектор. Таким образом, для определения оптимального управления $u=u(t)$ в задаче (10)-(12) нужно решить систему $(14),(15),(17),(18)$ относительно функций $u(t), \psi(t)$ и числа $\gamma \geqslant 0$. При $R=\infty$ здесь $\gamma=0, \psi(t)=0$ и условие (18) приведет к проблеме моментов (см. Моментов проблема): найти функцию $u=u(t)$, зная ее моменты



\[

\int_{t_{0}}^{T}\left\langle u(t), \quad \varphi_{k}(t)\right\rangle d t=a_{k}

\]



по системе $\varphi_{k}(t)=B^{T}(t) p_{k}(t), k=1, \ldots, n$.



Система (14), (15), (17), (18) представляет собой обобщенную проблему моментов для задачи (10)-(12) при $0<R<\infty$ (см. [3], [8]).



Јюбое решение $\psi(t)$ системы (15) однозначно представимо в виде



\[

\psi(t)=\sum_{k=1}^{n} \psi_{k} p_{k}(t), \quad t_{\theta} \leqslant t \leqslant T .

\]



При любом фиксированном $ү \geqslant 0$ существует решение $\psi(t ; \gamma), u(t ; \gamma)$ системы $(15),(17),(18)$, причем среди всех решений найдется единственное такое, что $u(t ; \gamma)$ имеет вид



\[

u(t ; \gamma)=\sum_{k=1}^{n} u_{k} B^{T}(t) p_{k}(t), \quad t_{0} \leqslant t \leqslant T .

\]



Для определения $\psi(t ; \gamma), u(t ; \gamma)$ нужно подставить выражения (19), (20) в (17), (18). В результате получится система линейных алгебраических уравнений относительно $\psi_{1}, \ldots, \psi_{n}, u_{1}, \ldots, u_{n}$, из к-рой однозначно определяются величины $\psi_{1}, \ldots, \psi_{n}$, а величины $u_{1}$, $\because \dot{\vec{T}}, u_{n}$ в случае линейной зависимости системы $\left\{\boldsymbol{B}^{\boldsymbol{T}}(t) p_{k}(t), k=1, \ldots, n\right\}$ определяются неоднозначно. сначала положить $\gamma=0, \psi(t)=0$ и из (18) определить $u(t ; 0)$ вида $(20)$. Затем следует проверить условие $\|u(t ; 0)\|_{L_{2}^{r}} \leqslant R$. Если это неравенство выполняется, то $u(t ; 0)-$ оптимальное управление задачи (10)-(12), имеющее минимальную норму среди всех оптимальных управлений; множество всех оптимальных управлений в этом случае исчерпывается управлениями вида



\[

u(t)=u(t ; 0)+v(t), \quad t_{0} \leqslant t \leqslant T,

\]



где $v(t)$ принадлежит ортогональному дополнению в $L_{2}^{r}\left[t_{0}, T\right]$ линейной оболочки систем функций



\[

\begin{aligned}

& \left\{B^{T}(t) p_{k}(t), \quad k=1, \ldots, n\right\}, \\

& \|v(t)\|_{L_{2}^{r}}^{2} \leqslant R^{2}-\|u(t ; 0)\|_{L_{2}^{r}}^{2} .

\end{aligned}

\]



Если $\| u\left(t ; 0 \|_{L_{2}^{r}}>R\right.$, то из (17), (18) при $\gamma>0$ ошределяют решения $\psi(t, \gamma), u(t, \gamma)$ вида $(19),(20)$ и находят $\gamma$ из уравнения



\[

\|u(t ; \gamma)\|_{L_{2}^{r}}=R, \quad \gamma>0

\]



функция $\|u(t ; \gamma)\|_{L_{2}^{r}}$ переменной $\gamma$ непрерывна, строго монотонно убывает при $\gamma>0, \lim _{\gamma \rightarrow+\infty} \| u\left(t ; 0 \|_{L_{2}^{r}}=0\right.$, поэтому из (21) однозначно определяется искомое $\gamma=\gamma_{*}$. Управление $u\left(t ; \gamma_{*}\right)$ будет оптимальным для задачи (10)-(12); при $\| u\left(t ; 0 \|_{L_{2}^{r}}>R\right.$ эта задача других оптимальных управлений не имеет.



Метод моментов применим также для решения вадачи быстродействия для систем (11) и друғих линейных систем (см. [3], [8]).



Упомянутые выше методы широко используются и для численного решения задач оптимального управления процессами, описываемыми уравнениями с частными производными.



Численная реализация многих методов рөшения задач оцтимального управления предполагает использование тех или иных методов приближенного решения





\end{document}